\section{Canonical Equivalence}
Any two regular Hamiltonian mechanical systems with the same number \(n\) of
degrees of freedom are locally related by a time-dependent canonical
transformation, on any interval where their Hamiltonian flows exist.

It is not generally true globally without additional assumptions. The cleanest
proof uses the Hamiltonian flow. A Hamilton–Jacobi proof gives the same result.

Let the two systems have canonical variables
\[
(q^i,p_i),\qquad (Q^i,P_i),\qquad i=1,\ldots,n
\]
and Hamiltonians
\[
H_1(q,p,t),\qquad H_2(Q,P,t).
\]
We claim that, locally in phase space and time, there exists a time-dependent
canonical transformation
\[
(q,p)\longleftrightarrow(Q,P)
\]
such that the transformed Hamiltonian of system 1 is precisely \(H_2\).
In other words, the two Hamiltonian systems are canonically equivalent.
The key point is that the explicit time dependence of the canonical
transformation is essential.

First step: any Hamiltonian system can be transformed into the trivial Hamiltonian
Consider one system
\[
H(q,p,t).
\]
Let
\[
\Phi_{t,t_0}:(q_0,p_0)\mapsto(q(t),p(t))
\]
be its Hamiltonian flow.
Because Hamiltonian flows are symplectic,
\[
\Phi_{t,t_0}^{*}\omega=\omega,
\]
where
\[
\omega=dq^i\wedge dp_i.
\]
Therefore the inverse flow
\[
\Phi_{t_0,t}
\]
is itself a canonical transformation.
Now define new variables by
\[
(Q,P)=\Phi_{t_0,t}(q,p).
\]
Thus \(Q,P\) are simply the initial phase-space coordinates of the trajectory
passing through \((q,p)\) at time \(t\).
Consequently,
\[
Q=\text{constant},\qquad P=\text{constant}
\]
along every Hamiltonian trajectory.
Hence
\[
\dot Q^i=0,\qquad \dot P_i=0.
\]
These are Hamilton's equations for the Hamiltonian
\[
\boxed{K(Q,P,t)=0}.
\]
Thus every Hamiltonian system can be transformed, by a time-dependent canonical
transformation, into
\[
\boxed{K=0}.
\]
This is the fundamental observation.

 More explicitly: why the transformed Hamiltonian is zero
For a time-dependent canonical transformation, the Hamiltonian does not
transform merely by substitution.
If the transformation is generated by a generating function \(F\), one has
schematically
\[
K=H+\frac{\partial F}{\partial t},
\]
with the precise form depending on the type of generating function.
The time dependence of the transformation therefore provides exactly the extra
term needed to cancel \(H\).
For example, take a type-2 generating function
\[
F_2(q,P,t).
\]
The transformation is
\[
p_i=\frac{\partial F_2}{\partial q^i},
\qquad
Q^i=\frac{\partial F_2}{\partial P_i},
\]
and
\[
K(Q,P,t)
=
H\left(q,\frac{\partial F_2}{\partial q},t\right)
+
\frac{\partial F_2}{\partial t}.
\]
If \(F_2\) satisfies
\[
\boxed{
\frac{\partial F_2}{\partial t}
+
H\left(q,\frac{\partial F_2}{\partial q},t\right)=0,
}
\]
then
\[
\boxed{K=0}.
\]
This is precisely the Hamilton–Jacobi equation.
A complete integral \(F_2(q,P,t)\), satisfying
\[
\det\left(
\frac{\partial^2F_2}
{\partial q^i\partial P_j}
\right)\neq0,
\]
defines the required local canonical transformation.

 Apply this to the first system
For
\[
H_1(q,p,t),
\]
choose a local complete solution of the Hamilton–Jacobi equation
\[
\frac{\partial S_1}{\partial t}
+
H_1\left(
q,\frac{\partial S_1}{\partial q},t
\right)=0.
\]
Define
\[
p_i=\frac{\partial S_1}{\partial q^i},
\]
and
\[
Q^i=\frac{\partial S_1}{\partial P_i}.
\]
Then
\[
K_1(Q,P,t)=0.
\]
Thus
\[
H_1
\quad\xrightarrow{\text{time-dependent canonical transformation}}\quad
0.
\]

 Do the same for the second system
For
\[
H_2(Q,P,t),
\]
choose a complete Hamilton–Jacobi solution \(S_2\). It gives a canonical transformation
\[
(Q,P)\longleftrightarrow (Q',P')
\]
such that
\[
H_2\longrightarrow K_2=0.
\]
Equivalently, reversing this transformation,
\[
0\longrightarrow H_2.
\]
Therefore we have the chain
\[
H_1
\longrightarrow
0
\longrightarrow
H_2.
\]
The composition of two canonical transformations is canonical.
Consequently,
\[
\boxed{
H_1
\stackrel{\rm can}{\longleftrightarrow}
H_2.
}
\]
This proves the claim locally.

An especially transparent formulation using flows
There is an even more geometric proof.
Let
\[
\Phi^{(1)}_t
\]
be the Hamiltonian flow of \(H_1\), and
\[
\Phi^{(2)}_t
\]
the Hamiltonian flow of \(H_2\).
Choose a reference time \(t_0\). Define
\[
\boxed{
T_t=
\Phi^{(2)}_{t,t_0}
\circ
\left(\Phi^{(1)}_{t,t_0}\right)^{-1}.
}
\]
Thus if \((q,p)\) is a point in the phase space of system 1 at time \(t\), we
first follow its trajectory backwards to \(t_0\), and then evolve that initial
point forward according to system 2.
Because both flows are symplectic,
\[
(\Phi^{(1)}_{t,t_0})^*\omega=\omega,
\]
and
\[
(\Phi^{(2)}_{t,t_0})^*\omega=\omega.
\]
Therefore
\[
\boxed{T_t^*\omega=\omega}.
\]
Hence \(T_t\) is a canonical transformation.
Now suppose system 1 follows
\[
x_1(t)=\Phi^{(1)}_{t,t_0}(x_0).
\]
Then
\[
T_t(x_1(t))
=
\Phi^{(2)}_{t,t_0}
\left[
(\Phi^{(1)}_{t,t_0})^{-1}
\Phi^{(1)}_{t,t_0}(x_0)
\right],
\]
so
\[
T_t(x_1(t))
=
\Phi^{(2)}_{t,t_0}(x_0).
\]
But the right-hand side is exactly a trajectory of system 2.
Therefore
\[
\boxed{
T_t:\text{ trajectories of system 1}
\longrightarrow
\text{trajectories of system 2}.
}
\]
This is perhaps the most direct proof of the proposition.

 Why the same number of degrees of freedom matters?
For \(n\) degrees of freedom, phase space has dimension
\[
\dim\Gamma=2n.
\]
The symplectic form is
\[
\omega=\sum_{i=1}^n dq^i\wedge dp_i.
\]
Two systems with different numbers of degrees of freedom have phase spaces of
different dimensions:
\[
2n_1\neq2n_2.
\]
There cannot be an invertible canonical transformation between them.
For equal \(n\), both phase spaces have the same symplectic dimension, and the
construction above can be made.
So the essential invariant here is the symplectic dimension.

 But there is an important caveat The statement should not be interpreted as
 saying that all mechanical systems are physically equivalent. The canonical
 transformation can be extremely complicated and explicitly time-dependent. For
 example, 
 \[ H_1=\frac{p^2}{2m}+\frac12m\omega^2q^2 \]
 and \[     H_2=\frac{P^2}{2M}+V(Q) \] 
 can locally be related in this sense, even if
 \(V(Q)\) produces chaotic motion while the oscillator is integrable. That does
 not mean that the oscillator and the chaotic system are equivalent in every
 physically relevant sense. The transformation itself may encode the complete
 dynamics of one system. Indeed, constructing it essentially requires knowing
 the Hamiltonian flow of the systems. So this theorem is a statement about
 canonical equivalence of Hamiltonian dynamics, not about physical equivalence
 under a simple change of variables.

Why the qualification "local" is important?
There are several possible obstructions to a global statement:
\begin{enumerate}
   \item Hamiltonian flows may not exist for all time.
   \item The Hamilton–Jacobi complete integral may exist only locally.
   \item The canonical transformation can become singular.
   \item The phase spaces may have different global topology.
   \item Constraints or singular Legendre transformations may invalidate the
       ordinary Hamiltonian description.
   \item Boundary conditions can distinguish systems globally.
   \item Global invariants can obstruct a single globally defined canonical map.
\end{enumerate}
For example, Darboux's theorem tells us that every symplectic manifold is locally
\[
(q^1,\ldots,q^n,p_1,\ldots,p_n)
\]
but says nothing about whether two symplectic manifolds are globally symplectomorphic.

. The precise theorem
A mathematically safe formulation is therefore:
Local canonical-equivalence theorem.
Let \(H_1\) and \(H_2\) be sufficiently regular time-dependent Hamiltonians
defined on \(2n\)-dimensional symplectic phase spaces, and suppose their
Hamiltonian flows exist on a common time interval. Then, locally in phase space
and on that interval, there exists a time-dependent canonical transformation
mapping the trajectories of \(H_1\) into those of \(H_2\).

The proof is simply
\[
\boxed{
T_t=
\Phi^{(2)}_{t,t_0}
\circ
(\Phi^{(1)}_{t,t_0})^{-1}.
}
\]
Since both factors are symplectic, \(T_t\) is canonical; since it maps one flow
into the other, it maps the corresponding Hamiltonian dynamics into each other.

A deeper point
This result is actually closely related to a rather striking fact about
Hamiltonian mechanics:
\emph{locally, the Hamiltonian itself carries no invariant information under
arbitrary time-dependent canonical transformations.}
The Hamiltonian can be transformed to zero.
What survives such transformations is not the particular function \(H(q,p,t)\),
but the symplectic structure together with the dynamical flow, including the
restrictions imposed on what transformations are regarded as physically
admissible.
This is also why, if one wants to compare two mechanical theories physically,
one normally restricts the allowed canonical transformations—for example to
time-independent transformations, transformations preserving a specified
configuration-space interpretation, locality, tensorial character, natural
Hamiltonian form, etc.
That restriction is where genuinely different mechanical theories can reappear.
